\providecommand{\SysName}{CARE}

\section{Conclusion}
\label{sec:conclusion}

\SysName{} demonstrates that security-aware refactoring for C/C++ can be framed
as a guarded synthesis problem: program analysis constructs the context needed
to identify high-value opportunities, an LLM proposes compact refactoring
patches, and a fail-closed validation loop filters candidates using build,
test, resource, semantic, and security-regression checks. Across the studied
projects, the results show that combining context construction with validation
is more reliable than using an LLM directly, especially for transformations
that touch cleanup paths, resource lifecycles, duplicated checks, and
error-handling behavior. While the current prototype is intentionally
conservative and does not provide formal equivalence guarantees, it establishes
a practical path toward refactoring tools that improve code structure without
weakening security-relevant behavior. More precise detectors, stronger semantic
checks, and structured edit generation can further improve \SysName{} and make
context-aware automated refactoring a useful component of secure software
maintenance.
